% !TITLE 梦李白·其一
% !AUTHOR 杜甫
% !DYNASTY 唐
% !ORDER 33

\begin{poem}
死别已吞声，生别常恻恻。
江南瘴疠地，逐客无消息。
故人入我梦，明我长相忆。
恐非平生魂，路远不可测。
魂来枫林青，魂返关塞黑。
君今在罗网，何以有羽翼？
落月满屋梁，犹疑照颜色。
水深波浪阔，无使蛟龙得。
\end{poem}

\begin{appreciation}
这是杜甫《梦李白》二首中的第一首。乾元二年（759年），李白因参与永王李璘幕府而获罪，被流放夜郎。杜甫不知道他已经在中途遇赦，日夜担心他的安危，终于在梦中见到了他。

死别已吞声，生别常恻恻——死别已经让人泣不成声，生别更是让人长久地悲伤。杜甫把生别和死别对比，说明生别比死别更折磨人：死别是一次的痛，生别是持续的痛。李白被流放，生死未卜，这种悬而未决的状态，是最让人煎熬的。

故人入我梦，明我长相忆——老朋友来到我的梦中，知道我长久地思念他。恐非平生魂，路远不可测——但恐怕这不是他生前的魂魄吧？路途那么远，谁知道发生了什么呢？魂来枫林青，魂返关塞黑——魂魄来的时候，江南的枫林正青；魂魄归去的时候，关塞一片漆黑。枫林青和关塞黑，一南一北，一青一黑，写出了空间的遥远和旅途的艰险。

君今在罗网，何以有羽翼——你现在身陷罗网，怎么还能有翅膀飞来呢？这是一个令人毛骨悚然的疑问：如果李白已经被囚禁，他的魂魄是怎么飞来的？杜甫不敢想答案，因为答案可能是他最不愿意面对的。

落月满屋梁，犹疑照颜色——梦醒之后，月光洒满屋梁，恍惚间还以为是照见了李白的面容。犹疑两个字，把梦醒后的恍惚写得极为真切。水深波浪阔，无使蛟龙得——江湖水深浪大，不要让蛟龙伤害了你。这是杜甫对李白的祝福，也是对自己无力保护朋友的愧疚。他只能在远方祈祷，希望李白平安。
\end{appreciation}
